% =====================================================================
%  ThmSmith Stage -1 proposal skeleton.
%  Written automatically by the ThmSmith TS pipeline before Stage -1 runs.
%  Locked parameters: qid=pid\_ordinal\_late\_partial\_defiers, specialization=v1
%
%  HOW TO USE THIS FILE
%  --------------------
%  - The preamble (everything above the `% --- BEGIN CONTENT ---` marker)
%    and the `\section{...}` headers below are fixed by the pipeline.
%    Do NOT rewrite or rename them; the Step 3 reviewer subagent and the
%    downstream Stage 0 / Stage 1 prompts grep for these section titles.
%  - Each section contains exactly one `% TODO(stage-1 ...): ...`
%    placeholder. Replace each TODO with the content described for that
%    section in `ThmSmith/tools/src/prompts/stage_neg1_draft.txt`
%    ("REQUIRED OUTPUT FORMAT (Stage -1 proposal .tex)").
%  - The `\paragraph{Locked parameters.}` block in the metadata block
%    must pin the chosen `(qid, specialization, p, assignment, target,
%    regression)` precisely. If the proposal maps to a Q1..Q6 pool-mode,
%    reproduce that block verbatim from
%    `ThmSmith/doc/panel_formula_question_generator_note_with_common_prompts.tex`.
%    If it is a free-form `(qid, specialization)`, define each parameter
%    explicitly here (no implicit conventions). Stage 0 quotes this block;
%    do not rephrase it after locking.
%  - Removing a section header, renaming a macro, or rewriting the
%    preamble is a regression: the file must still match this skeleton
%    after you finish.
% =====================================================================

\documentclass[11pt]{article}

\usepackage{amsmath,amssymb,amsthm,mathtools}
\usepackage{enumitem}
\usepackage[colorlinks=true,linkcolor=blue,citecolor=blue,urlcolor=blue]{hyperref}
\usepackage{natbib}

\newcommand{\E}{\mathbb{E}}
\renewcommand{\Pr}{\mathbb{P}}
\newcommand{\one}{\mathbf{1}}
\newcommand{\transpose}{\mathsf{T}}
\newcommand{\calH}{\mathcal{H}}
\newcommand{\calF}{\mathcal{F}}
\newcommand{\calR}{\mathcal{R}}
\newcommand{\R}{\mathbb{R}}
\newcommand{\plim}{\operatorname*{plim}}

\theoremstyle{definition}
\newtheorem{definition}{Definition}
\newtheorem{assumption}{Assumption}
\newtheorem{example}{Example}

\theoremstyle{plain}
\newtheorem{theorem}{Theorem}
\newtheorem{conjecture}{Conjecture}
\newtheorem{proposition}{Proposition}
\newtheorem{lemma}{Lemma}
\newtheorem{corollary}{Corollary}

\theoremstyle{remark}
\newtheorem{remark}{Remark}

\title{ThmSmith Stage -1 proposal: \texttt{pid\_ordinal\_late\_partial\_defiers} / \texttt{v1}%
  \\ \large\textit{Zig-zag compatibility tightens ordinal-IV LATE bounds under per-rung defier caps}}
\author{ThmSmith Stage -1 proposer}
\date{\today}

\begin{document}
\maketitle

% --- BEGIN CONTENT ---  (everything above is fixed by the pipeline)

\paragraph{Locked parameters.}
This is a partial-identification anchor, not a panel pool-mode.  The locked tuple is
\[
\begin{gathered}
(\texttt{qid},\texttt{specialization},P,\theta,\text{identifying restrictions},
  \text{bounding functional})\\
=
(\texttt{pid\_ordinal\_late\_partial\_defiers},\texttt{v1},
  P^{(1)},\theta^+_\omega,\mathcal R_{\eta}^{\mathrm{ordIV}},
  [\underline\theta_\omega(P,\eta),\overline\theta_\omega(P,\eta)]) .
\end{gathered}
\]
Here \(P^{(1)}=\{p_z,d_z,m_z:z=0,\ldots,J\}\) is the observable first-moment array with \(p_z=\Pr(Z=z)>0\), \(d_z=\E[D\mid Z=z]\), and \(m_z=\E[Y\mid Z=z]\).  The outcome is bounded, \(Y\in[0,1]\), \(Z\in\{0,\ldots,J\}\) is an ordinal instrument with \(J\ge 1\), \(D\in\{0,1\}\), and the full-data response path is \(R=(D(0),\ldots,D(J))\in\{0,1\}^{J+1}\).  For adjacent rung \(j\), let
\[
U_j=\{D(j)=0,D(j+1)=1\},\qquad V_j=\{D(j)=1,D(j+1)=0\}
\]
be the up-complier and down-defier events, with masses \(q^+_j=\Pr(U_j)\) and \(q^-_j=\Pr(V_j)\).  The target is the weighted positive-rung LATE
\[
\theta^+_\omega
=\frac{\sum_{j=0}^{J-1}\omega_j\,\E[(Y(1)-Y(0))\one\{U_j\}]}
        {\sum_{j=0}^{J-1}\omega_j q^+_j},
\]
where the known weights satisfy \(\omega_j>0\).  The restrictions \(\mathcal R_{\eta}^{\mathrm{ordIV}}\) are ordinal-IV random assignment, exclusion, consistency, bounded outcomes, positive instrument cells, \(q^-_j\le \eta_j\) for known per-rung defier caps \(\eta_j\), and the observable feasibility inequalities \(\max\{0,d_j-d_{j+1}\}\le\eta_j\).  The bounding functional is the sharp projection of the finite response-path latent-type polytope matching \(P^{(1)}\), with the pairwise-pasting relaxation defined in Assumption~\ref{ass:pairwise}.

\begin{abstract}
This proposal asks for sharp partial-identification bounds on an ordinal-instrument LATE when Imbens--Angrist monotonicity is relaxed to allow a calibrated fraction of down-defiers on each adjacent instrument rung.  The routine result is a finite response-path sharpness theorem: the identified set is the projection of one latent table over full ordinal treatment-response paths, not a product of adjacent binary-IV tables.  The novel kernel is the zig-zag compatibility conjecture: with three instrument levels, adjacent pairwise bounded-defier LATE bounds can be strictly too wide because the same response path \(010\) is an up-complier on one rung and a down-defier on the next, so its treatment effect must be shared with opposite signs across reduced forms.  This phenomenon is invisible in binary-instrument sensitivity analyses and in full-monotonicity ordinal-IV formulas.  If resolved, it would give a formal criterion for when per-rung monotonicity sensitivity can be calibrated locally and when the ordinal response-path geometry itself tightens the LATE bounds.
\end{abstract}

\section{Introduction}
\paragraph{Why the problem is important.}
Ordinal instruments are common when encouragement or eligibility intensity has several ordered levels.  Under full Imbens--Angrist monotonicity, adjacent reduced forms identify a nonnegative weighted average of local causal effects.  Applied researchers often find full no-defiers monotonicity too strong, but a binary sensitivity analysis treats each adjacent rung in isolation and misses the fact that the same unit has one whole response path across all rungs.

\paragraph{The main finding (proposed).}
Theorem~\ref{thm:monotone-reduction} records the monotone benchmark: when all \(\eta_j=0\), the target reduces to the usual ordinal adjacent-rung weighted Wald estimand.  Theorem~\ref{thm:finite-polytope} gives the finite sharpness representation under bounded per-rung defier proportions.  Conjecture~\ref{conj:zigzag} is the flagship novelty claim: for \(J=2\), the sharp global response-path bound can be strictly tighter than the natural pairwise-pasting relaxation formed by solving two adjacent binary-IV bounded-defier problems separately.

\paragraph{What is new.}
The closest failed in-repo angle, \texttt{pid\_late\_bounded\_defiers}/\texttt{one\_sided\_compliance}, studied a binary-instrument bounded-defier envelope.  The present proposal is not a re-run of that kernel: its new mechanism requires at least three instrument rungs and a response path such as \(010\) or \(101\).  Such a path contributes to two adjacent reduced forms with opposite signs, so independently sharp adjacent bounds can require incompatible values of the same latent treatment effect.  The novelty is the cross-rung compatibility facet, not the existence of a bounded-defier binary-IV interval.

\paragraph{How it positions in the literature.}
This proposal lives in the partial-identification cluster for IV/LATE bounds and monotonicity sensitivity.  It is downstream of Imbens--Angrist ordinal-IV monotonicity and Balke--Pearl latent-type sharp bounds, adjacent to LATE-with-defiers work by de Chaisemartin and by Dahl--Huber--Mellace, and closest in pairwise spirit to Sun--Wuthrich pairwise valid instruments.  It is also complementary to recent results showing limits of monotonicity's identifying power for ATE-type targets.  It deliberately avoids the panel cluster and does not propose new SCM or do-calculus infrastructure.

\section{Related work}
\citet{ImbensAngrist1994} establish binary-instrument LATE identification under independence, exclusion, and monotonicity; this proposal relaxes the no-defiers clause.  \citet{AngristImbens1995} study variable treatment intensity and show that IV estimands can be interpreted as weighted averages of causal responses under ordered monotonicity; the proposed kernel asks what remains sharp when adjacent ordered monotonicity is only approximately true.  \citet{AngristImbensRubin1996} give the principal-strata formulation that makes the up-complier and down-defier events precise.  \citet{Vytlacil2002} links independence and monotonicity to latent-index selection; zig-zag response paths are exactly the ordinal-pattern obstruction to that scalar-index picture.  \citet{BalkePearl1997} derive sharp IV bounds by optimizing over latent response types, motivating the finite response-path polytope in Theorem~\ref{thm:finite-polytope}.

\citet{deChaisemartin2017} shows that IV can retain a LATE interpretation with defiers under a compliers-defiers condition; this proposal instead keeps only per-rung defier caps and accepts partial identification.  \citet{HuberMellace2012} replace global monotonicity by local monotonicity that separates compliers and defiers through outcome-support regions; this is a close binary-IV local-monotonicity comparison, but it is not an ordinal response-path sharp-bounds result and does not analyze pairwise-versus-global non-pasteability.  \citet{DahlHuberMellace2023} identify LATEs for compliers and defiers under local monotonicity and local compliers-defiers assumptions; the present bounded-defier model is weaker and should produce a sharp interval with cross-rung facets.  \citet{Noack2021} studies sensitivity of LATE estimates to monotonicity violations using defier-share and outcome-distribution parameters; this proposal isolates the additional ordinal compatibility created by a single response path across several adjacent pairs.  \citet{SunWuthrich2025} introduce pairwise valid instruments and estimate LATEs for researcher-selected valid instrument-value pairs with researcher-specified aggregation weights; the comparator is close because it is pairwise, but it treats valid pairs as local estimands rather than deriving a sharp global response-path interval under bounded adjacent defier caps.  \citet{MogstadTorgovitskyWalters2021} show that multiple-IV 2SLS needs partial monotonicity conditions for a positive-weight LATE interpretation; this proposal studies a single ordinal instrument but a related local monotonicity relaxation.  \citet{BaiHuangMoonShaikhVytlacil2024} show that generalized monotonicity often has no identifying power for ATE sets; the contrast here is a latent-type LATE target whose denominator and signs depend directly on the response-path pattern.  \citet{HuberMellace2015} and \citet{Kitagawa2015} provide inequality-based tests of LATE assumptions, while this proposal turns adjacent monotonicity slack into quantitative sharp bounds.

In-repo prior art is sparse for this exact anchor.  \texttt{doc/API.md} section 3 records \texttt{ThmSmith/PartialID/} as the staging area for Manski-family bounds, Balke--Pearl, IV bounds, sensitivity analysis, and shape restrictions, but no current ThmSmith catalogue entry has an ordinal-IV bounded-defier response-path theorem.  The closest banked proposal is \texttt{pid\_late\_bounded\_defiers}/\texttt{one\_sided\_compliance}, whose headline is a binary-instrument one-sided closed-form envelope; it lacks ordinal rungs and therefore cannot contain the zig-zag compatibility phenomenon.  The accepted banked \texttt{pid\_dynamic\_iv\_compliance\_v2} proposal uses full-history dynamic IV bounds; it is useful as a latent response-map sharpness analogue, but its kernel is assignment-memory retention rather than adjacent ordinal monotonicity sensitivity.

\section{Formal setup}
Let \(P=P^{(1)}\) be the observable first-moment array \(\{p_z,d_z,m_z:z=0,\ldots,J\}\), with \(p_z=\Pr(Z=z)>0\), \(d_z=\E[D\mid Z=z]\), and \(m_z=\E[Y\mid Z=z]\).  The instrument satisfies \(Z\in\{0,\ldots,J\}\), \(J\ge1\), the treatment is binary \(D\in\{0,1\}\), and outcomes obey \(Y\in[0,1]\).  A full-data law contains potential treatments \(D(z)\) for all instrument levels and potential outcomes \(Y(0),Y(1)\).  Exclusion and consistency give \(D=D(Z)\) and \(Y=Y(D)\).  The response path is \(R=(D(0),\ldots,D(J))\).

For each adjacent rung \(j\in\{0,\ldots,J-1\}\), define the up-complier event \(U_j=\{D(j)=0,D(j+1)=1\}\) and the down-defier event \(V_j=\{D(j)=1,D(j+1)=0\}\).  The observed first stage is
\[
\Delta^D_j=d_{j+1}-d_j=\E[D\mid Z=j+1]-\E[D\mid Z=j]=q^+_j-q^-_j,
\]
and the observed reduced form is
\[
\Delta^Y_j=m_{j+1}-m_j=\E[Y\mid Z=j+1]-\E[Y\mid Z=j]
=\E[(Y(1)-Y(0))\one\{U_j\}]
-\E[(Y(1)-Y(0))\one\{V_j\}].
\]
The target is \(\theta^+_\omega\), the weighted average treatment effect over positive adjacent response events \(U_j\), with weights \(\omega_j>0\).  The identified set is
\[
\Theta_I(P,\eta)
=\{\theta^+_\omega(Q): Q \text{ is a full-data law whose induced first moments equal } P
   \text{ and satisfies } \mathcal R_\eta^{\mathrm{ordIV}}\}.
\]
The sharp bounding functional is
\[
[\underline\theta_\omega(P,\eta),\overline\theta_\omega(P,\eta)]
=[\inf\Theta_I(P,\eta),\sup\Theta_I(P,\eta)].
\]

\begin{verbatim}
qid = pid_ordinal_late_partial_defiers
specialization = v1
observable distribution P = first-moment array {p_z,d_z,m_z}, where
  p_z = Pr(Z=z), d_z = E[D | Z=z], and m_z = E[Y | Z=z]
parameter theta = theta^+_omega, the weighted positive-adjacent-rung LATE
identifying restrictions = ordinal IV independence/exclusion, consistency,
  bounded outcomes, positive instrument cells, known per-rung defier caps eta_j
bounding functional = sharp projection of the finite response-path latent-type
  polytope matching the first-moment array; pairwise-pasting relaxation used
  only as a comparator
\end{verbatim}

\section{Assumptions}
\begin{assumption}[Ordinal finite IV setup]\label{ass:finite-ordinal}
\(Z\in\{0,\ldots,J\}\) with \(J\ge1\), \(D\in\{0,1\}\), \(Y\in[0,1]\), and the known target weights satisfy \(\omega_j>0\) for every adjacent rung \(j\).
\end{assumption}

\begin{assumption}[Ordinal IV validity]\label{ass:ord-iv}
\((D(0),\ldots,D(J),Y(0),Y(1))\) is independent of \(Z\), \(D=D(Z)\), and \(Y=Y(D)\) almost surely.
\end{assumption}

\begin{assumption}[Positive cells and positive target mass]\label{ass:positive}
\(\Pr(Z=z)>0\) for every \(z\in\{0,\ldots,J\}\), and \(\sum_{j=0}^{J-1}\omega_j q^+_j>0\).
\end{assumption}

\begin{assumption}[Known per-rung defier caps]\label{ass:defier-caps}
For known constants \(\eta_j\in[0,1]\), the adjacent down-defier masses satisfy \(q^-_j=\Pr(V_j)\le\eta_j\).  The observable first-stage feasibility condition \(\max\{0,-\Delta^D_j\}\le\eta_j\) holds for every \(j\).
\end{assumption}

\begin{assumption}[Finite response-path sharpness class]\label{ass:finite-types}
Sharpness is evaluated over all response paths \(r\in\{0,1\}^{J+1}\).  Let \(\pi_r\) be the path mass and let \(\mu_{0r},\mu_{1r}\) be outcome-mass variables, where \(\mu_{dr}=\pi_r y_d(r)\) for some path-specific potential-outcome mean \(y_d(r)\in[0,1]\).  Equivalently, the endpoint program uses the linear constraints \(0\le\mu_{dr}\le\pi_r\), \(\sum_r\pi_r=1\), and \(\pi_r\ge0\).  The induced moments must equal \(P\): for every \(z\),
\[
\sum_r \pi_r r_z=d_z,\qquad
\sum_r \{r_z\mu_{1r}+(1-r_z)\mu_{0r}\}=m_z.
\]
The class must also satisfy Assumptions~\ref{ass:finite-ordinal}--\ref{ass:defier-caps}; in particular, each cap is \(\sum_{r:r_j=1,r_{j+1}=0}\pi_r\le\eta_j\).
\end{assumption}

\begin{assumption}[Pairwise-pasting comparator]\label{ass:pairwise}
For each adjacent pair \(j\), let \(\mathcal F_j(P,\eta_j)\) be the binary-IV bounded-defier feasible set for the two cells \((Z=j,Z=j+1)\), matching \((d_j,m_j,d_{j+1},m_{j+1})\) and satisfying \(q^-_j\le\eta_j\).  Each element of \(\mathcal F_j\) supplies an adjacent up-complier numerator \(a_j=\E[(Y(1)-Y(0))\one\{U_j\}]\) and denominator \(b_j=q^+_j\).  The pairwise-pasting interval is
\[
B_{\mathrm{pair}}(P,\eta)=
\left[
\inf_{(a_j,b_j)\in\mathcal F_j\ \forall j}
\frac{\sum_j\omega_j a_j}{\sum_j\omega_j b_j},
\;
\sup_{(a_j,b_j)\in\mathcal F_j\ \forall j}
\frac{\sum_j\omega_j a_j}{\sum_j\omega_j b_j}
\right],
\]
where the product relaxation optimizes the adjacent binary problems independently before forming the aggregate ratio, and the denominator is restricted to be positive.
\end{assumption}

\begin{assumption}[Three-rung zig-zag witness class]\label{ass:zigzag-class}
For Conjectures~\ref{conj:zigzag} and~\ref{conj:kink}, take \(J=2\), \(\omega_0=\omega_1=1\), positive first stages \(\Delta^D_0,\Delta^D_1>0\), and caps satisfying \(0<\eta_0,\eta_1<1\).  The response paths \(010\) and \(101\) have positive admissible mass under the response-path polytope.  For Conjecture~\ref{conj:kink}, cap sensitivity is evaluated only along a fixed positive ray \(\eta(t)=t v\), where \(v=(v_0,v_1)\in(0,\infty)^2\) is fixed and \(t\ge0\).
\end{assumption}

\section{Main results}
\begin{theorem}[Theorem 1: monotone ordinal-IV reduction]\label{thm:monotone-reduction}
Under Assumptions~\ref{ass:finite-ordinal}--\ref{ass:defier-caps}, if \(\eta_j=0\) for every \(j\), then
\[
\theta^+_\omega
=
\frac{\sum_{j=0}^{J-1}\omega_j\Delta^Y_j}
     {\sum_{j=0}^{J-1}\omega_j\Delta^D_j}.
\]
Equivalently, the sharp identified set is a singleton equal to the adjacent-rung weighted Wald estimand.
\end{theorem}
\paragraph{Why this is new and worth studying.}
This is not the novelty kernel; it is the sanity anchor to \citet{ImbensAngrist1994} and \citet{AngristImbens1995}.  It ensures the bounded-defier proposal reduces to the standard ordered-monotonicity estimand rather than inventing a new target by notation.

\begin{theorem}[Theorem 2: finite response-path sharpness]\label{thm:finite-polytope}
Under Assumptions~\ref{ass:finite-ordinal}--\ref{ass:finite-types}, the sharp identified set for \(\theta^+_\omega\) relative to the first-moment observable array \(P\) is the projection of one finite latent response-path feasible set.  More precisely, the lower and upper endpoints \(\underline\theta_\omega(P,\eta)\) and \(\overline\theta_\omega(P,\eta)\) are attained by a finite linear-fractional program over variables \((\pi_r,\mu_{0r},\mu_{1r})_{r\in\{0,1\}^{J+1}}\) satisfying the linear constraints in Assumption~\ref{ass:finite-types}.  The objective is
\[
\frac{\sum_{j=0}^{J-1}\omega_j
  \sum_{r:r_j=0,r_{j+1}=1}(\mu_{1r}-\mu_{0r})}
{\sum_{j=0}^{J-1}\omega_j
  \sum_{r:r_j=0,r_{j+1}=1}\pi_r}.
\]
Since Assumption~\ref{ass:positive} makes the target denominator positive, the endpoint problems are equivalent to finite linear programs by the standard Charnes--Cooper rescaling.
\end{theorem}
\paragraph{Why this is new and worth studying.}
The closest prior art is \citet{BalkePearl1997}, plus the in-repo Balke--Pearl and binary bounded-defier proposal families.  The difference is the full ordinal response-path index: adjacent monotonicity violations are not separate binary response types but coordinates of one path \(R\), which is exactly what later conjectures exploit.

\begin{conjecture}[Conjecture 1: zig-zag compatibility strictly tightens pairwise bounds]\label{conj:zigzag}
(NOT YET PROVED.)  Under Assumptions~\ref{ass:finite-ordinal}--\ref{ass:pairwise} and the witness class in Assumption~\ref{ass:zigzag-class}, there exist an observable first-moment array \(P\) and caps \(\eta=(\eta_0,\eta_1)\) such that
\[
[\underline\theta_\omega(P,\eta),\overline\theta_\omega(P,\eta)]
\subsetneq B_{\mathrm{pair}}(P,\eta).
\]
The strict inclusion is generated by zig-zag paths: a unit with response path \(010\) contributes positively to the \(0\to1\) reduced form and negatively to the \(1\to2\) reduced form, while the path \(101\) contributes with the opposite sign pattern.  Pairwise binary-IV relaxations can assign incompatible potential-outcome effects to these two appearances; the global ordinal response-path feasible set cannot.
\end{conjecture}
\paragraph{Why this is new and worth studying.}
\citet{HuberMellace2012} is a close defier reference because it shows how local monotonicity can separate compliers and defiers by outcome-support regions, but it is a binary-IV local-monotonicity identification result rather than an ordinal bounded-defier sharp-bounds result.  \citet{SunWuthrich2025} is the closest pairwise comparator because it estimates LATEs for valid instrument-value pairs and aggregates them with researcher-specified weights, but it does not impose one global ordinal response path or ask whether independently valid pairwise intervals can be pasted into a sharp global interval.  \citet{Noack2021} and the failed in-repo binary bounded-defier proposal study monotonicity sensitivity for a binary instrument, where no cross-rung path can exist.  \citet{deChaisemartin2017} and \citet{DahlHuberMellace2023} restore point identification under additional cancellation or local monotonicity conditions, whereas this conjecture predicts a partial-ID tightening under only per-rung caps.  Proving it would tell researchers when adjacent-rung sensitivity analyses are conservative because they ignore ordinal response-path compatibility.

\begin{conjecture}[Conjecture 2: first ordinal kink at the zig-zag facet]\label{conj:kink}
(NOT YET PROVED.)  In the \(J=2\) witness class of Assumption~\ref{ass:zigzag-class}, fix a positive ray \(v\in(0,\infty)^2\) and set \(\eta(t)=t v\).  The one-dimensional endpoint functions \(\underline\theta_\omega(P,\eta(t))\) and \(\overline\theta_\omega(P,\eta(t))\) are piecewise linear-fractional in \(t\).  There are observable first-moment arrays \(P\) and rays \(v\) for which the first breakpoint \(t^\star>0\) from the monotone point \(t=0\) is not a single-rung binary-IV endpoint but the ordinal compatibility facet that limits the joint mass and common treatment-effect assignment of the zig-zag paths \(010\) and \(101\).
\end{conjecture}
\paragraph{Why this is new and worth studying.}
\citet{HuberMellace2012} and \citet{DahlHuberMellace2023} show how local monotonicity can recover identifiable complier and defier parameters, but they do not analyze the endpoint geometry of calibrated defier shares.  \citet{SunWuthrich2025} is pairwise in a different sense: it selects valid pairs and estimates their LATEs, while this conjecture studies the first parametric sensitivity breakpoint caused by incompatibility across adjacent pairs that share response paths.  \citet{MogstadTorgovitskyWalters2021} and \citet{BaiHuangMoonShaikhVytlacil2024} clarify what monotonicity-type restrictions buy for IV estimands, but not the geometry of per-rung sensitivity bounds for an ordinal response path.  A ray-indexed kink theorem would give a concise diagnostic: if the active endpoint facet is zig-zag compatibility, then the ordinal design supplies information that no adjacent binary analysis can see.

\section{Sanity-check corollaries}
\begin{corollary}[Binary-instrument reduction]\label{cor:binary}
Under Assumptions~\ref{ass:finite-ordinal}--\ref{ass:finite-types}, if \(J=1\), then the response-path polytope in Theorem~\ref{thm:finite-polytope} has only the four binary principal strata \(00,01,10,11\), and \(B_{\mathrm{pair}}(P,\eta)\) equals the global sharp bound.
\end{corollary}
Under Assumptions~\ref{ass:finite-ordinal}--\ref{ass:finite-types}, the reduction is exact because there is only one adjacent rung, so the product in Assumption~\ref{ass:pairwise} has one factor and a response path cannot be a complier for one rung and a defier for another.

\begin{corollary}[Full-monotonicity reduction]\label{cor:full-monotone}
Under Assumptions~\ref{ass:finite-ordinal}--\ref{ass:defier-caps}, if \(\eta_j=0\) for every \(j\), all down-defier events \(V_j\) have zero mass and Theorem~\ref{thm:monotone-reduction} recovers the ordered adjacent Wald estimand of \citet{AngristImbens1995}.
\end{corollary}
Under Assumptions~\ref{ass:finite-ordinal}--\ref{ass:defier-caps}, the reduction is algebraic: \(\eta_j=0\) and \(q^-_j\le\eta_j\) imply \(q^-_j=0\), hence \(\Delta^Y_j=\E[(Y(1)-Y(0))\one\{U_j\}]\) and \(\Delta^D_j=q^+_j\) on every rung.

\section{Proof-sketch route}
The verification route is finite-dimensional.  Expand the ordinal IV model into response paths \(r\in\{0,1\}^{J+1}\), write each observed treatment mean \(d_z\) as a linear functional of \(\pi_r\), write each observed outcome mean \(m_z\) as a linear functional of outcome masses \(\mu_{0r},\mu_{1r}\), and impose \(q^-_j\le\eta_j\) as linear inequalities in \(\pi_r\).  The monotone reduction follows from \(q^-_j=0\) after combining \(\eta_j=0\) with the cap assumption.  For sharpness, optimize over the resulting bounded polytope, then convert the ratio target to a linear program by positive-denominator rescaling.  The conjectural part is a \(J=2\) small-support enumeration: compare the global eight-path feasible set against two independently optimized four-type adjacent relaxations, and exhibit a rational observable law whose endpoint in the pairwise relaxation requires inconsistent \(\mu_{1r}-\mu_{0r}\) assignments for the shared \(010\) or \(101\) path.  The smallest witness search should enumerate rational grids for \((d_0,d_1,d_2,m_0,m_1,m_2,\eta_0,\eta_1)\), solve the global and pairwise linear-fractional endpoints after Charnes--Cooper rescaling, and record the first strict-inclusion certificate by the active constraints.  The kink claim should be checked along a fixed ray \(\eta(t)=t v\), so the active-facet sequence is a one-dimensional parametric linear-program calculation rather than a path-dependent two-dimensional statement.

\section{Honest scope flags}
Theorem~\ref{thm:monotone-reduction} and Theorem~\ref{thm:finite-polytope} are labelled as theorems because they are routine finite response-type algebra.  Conjecture~\ref{conj:zigzag} and Conjecture~\ref{conj:kink} are the novelty kernel and are not yet proved.  The proposal does not claim that per-rung caps are point-identifying, and it does not claim that every ordinal-IV sensitivity analysis is tighter than its pairwise relaxation; strict tightening is asserted only as an existence phenomenon in the \(J=2\) witness class.  A Stage 0 derivation should either produce a rational \(J=2\) witness or downgrade Conjecture~\ref{conj:kink} to a search problem.

\section{Iteration trail}
\begin{itemize}[leftmargin=*]
\item v1 (cold-start) -- initial draft proposing ordinal response-path zig-zag compatibility as the novelty kernel; prior verdict: none.
\item v2 (revise) -- added Huber--Mellace 2012 as closest ordinal-defier prior art, repaired the defier-cap hypothesis in Theorem~\ref{thm:monotone-reduction}, changed \(P\) to a first-moment observable array for the finite program, made \(B_{\mathrm{pair}}\) an explicit aggregate-ratio product relaxation, and ray-indexed the kink conjecture; prior verdict: REVISE.
\item v3 (revise) -- corrected the Huber--Mellace comparison, added Sun--Wuthrich 2025 as the pairwise-valid-instrument comparator, and restated Theorem~\ref{thm:finite-polytope} with outcome-mass variables \(\mu_{dr}\) so the endpoint program is genuinely linear-fractional; prior verdict: REVISE.
\end{itemize}

\section*{Bibliography}
\begin{thebibliography}{99}
\bibitem[Angrist and Imbens(1995)]{AngristImbens1995}
Angrist, J. D. and G. W. Imbens (1995).
\newblock Two-stage least squares estimation of average causal effects in models with variable treatment intensity.
\newblock \emph{Journal of the American Statistical Association} 90(430), 431--442.
\newblock \url{https://doi.org/10.1080/01621459.1995.10476535}.

\bibitem[Angrist, Imbens, and Rubin(1996)]{AngristImbensRubin1996}
Angrist, J. D., G. W. Imbens, and D. B. Rubin (1996).
\newblock Identification of causal effects using instrumental variables.
\newblock \emph{Journal of the American Statistical Association} 91(434), 444--455.
\newblock \url{https://doi.org/10.1080/01621459.1996.10476902}.

\bibitem[Bai et~al.(2024)]{BaiHuangMoonShaikhVytlacil2024}
Bai, Y., S. Huang, S. Moon, A. M. Shaikh, and E. J. Vytlacil (2024).
\newblock On the identifying power of generalized monotonicity for average treatment effects.
\newblock NBER Working Paper 32983.
\newblock \url{https://doi.org/10.3386/w32983}.

\bibitem[Balke and Pearl(1997)]{BalkePearl1997}
Balke, A. and J. Pearl (1997).
\newblock Bounds on treatment effects from studies with imperfect compliance.
\newblock \emph{Journal of the American Statistical Association} 92(439), 1171--1176.
\newblock \url{https://doi.org/10.1080/01621459.1997.10474074}.

\bibitem[Dahl, Huber, and Mellace(2023)]{DahlHuberMellace2023}
Dahl, C. M., M. Huber, and G. Mellace (2023).
\newblock It is never too LATE: a new look at local average treatment effects with or without defiers.
\newblock \emph{The Econometrics Journal} 26(3), 378--404.
\newblock \url{https://doi.org/10.1093/ectj/utad013}.

\bibitem[de Chaisemartin(2017)]{deChaisemartin2017}
de Chaisemartin, C. (2017).
\newblock Tolerating defiance? Local average treatment effects without monotonicity.
\newblock \emph{Quantitative Economics} 8(2), 367--396.
\newblock \url{https://doi.org/10.3982/QE601}.

\bibitem[Huber and Mellace(2015)]{HuberMellace2015}
Huber, M. and G. Mellace (2015).
\newblock Testing instrument validity for LATE identification based on inequality moment constraints.
\newblock \emph{Review of Economics and Statistics} 97(2), 398--411.
\newblock \url{https://doi.org/10.1162/REST_a_00450}.

\bibitem[Huber and Mellace(2012)]{HuberMellace2012}
Huber, M. and G. Mellace (2012).
\newblock Relaxing monotonicity in the identification of local average treatment effects.
\newblock Economics Working Paper Series 1212, University of St. Gallen, School of Economics and Political Science.
\newblock \url{https://ux-tauri.unisg.ch/RePEc/usg/econwp/EWP-1212.pdf}.

\bibitem[Imbens and Angrist(1994)]{ImbensAngrist1994}
Imbens, G. W. and J. D. Angrist (1994).
\newblock Identification and estimation of local average treatment effects.
\newblock \emph{Econometrica} 62(2), 467--475.
\newblock \url{https://www.nber.org/papers/t0118}.

\bibitem[Kitagawa(2015)]{Kitagawa2015}
Kitagawa, T. (2015).
\newblock A test for instrument validity.
\newblock \emph{Econometrica} 83(5), 2043--2063.
\newblock \url{https://doi.org/10.3982/ECTA11974}.

\bibitem[Mogstad, Torgovitsky, and Walters(2021)]{MogstadTorgovitskyWalters2021}
Mogstad, M., A. Torgovitsky, and C. R. Walters (2021).
\newblock The causal interpretation of two-stage least squares with multiple instrumental variables.
\newblock \emph{American Economic Review} 111(11), 3663--3698.
\newblock \url{https://doi.org/10.1257/aer.20190221}.

\bibitem[Noack(2021)]{Noack2021}
Noack, C. (2021).
\newblock Sensitivity of LATE estimates to violations of the monotonicity assumption.
\newblock arXiv:2106.06421.
\newblock \url{https://arxiv.org/abs/2106.06421}.

\bibitem[Sun and W\"uthrich(2025)]{SunWuthrich2025}
Sun, Z. and K. W\"uthrich (2025).
\newblock Pairwise valid instruments.
\newblock \emph{Journal of Econometrics} 250, 106009.
\newblock \url{https://doi.org/10.1016/j.jeconom.2025.106009}.

\bibitem[Vytlacil(2002)]{Vytlacil2002}
Vytlacil, E. (2002).
\newblock Independence, monotonicity, and latent index models: An equivalence result.
\newblock \emph{Econometrica} 70(1), 331--341.
\newblock \url{https://doi.org/10.1111/1468-0262.00282}.
\end{thebibliography}

\end{document}
